\section{机器人数据采集方式：真实、遥操作、仿真、合成}

机器人数据采集有多种路径，每一种都有成本和偏差。真实部署最贴近产品，但规模不足；遥操作能得到动作示范，但人力成本高；仿真能规模化长尾场景，但有 sim-to-real gap；互联网视频规模大，但缺少动作和反馈。高标准数据系统必须组合这些来源。

\subsection{采集方式对比}

\noindent\begin{longtable}{p{0.22\textwidth}p{0.30\textwidth}p{0.39\textwidth}}
\toprule
方式 & 优点 & 代价 \\
\midrule
真实部署 & 最贴近真实任务和硬件 & 规模小、风险高、长尾少。 \\
遥操作 & 动作示范清晰，可控性强 & 人力昂贵，操作者分布有限。 \\
仿真 & 可规模生成、可覆盖危险/稀有场景 & 真实差距、物理建模成本。 \\
合成数据 & 可快速补齐任务和语言标注 & 分布偏差和幻觉风险。 \\
互联网视频 & 规模巨大，覆盖广泛场景 & 缺动作、缺状态、缺 reward。 \\
\bottomrule
\end{longtable}

\begin{knowledgebox}{为什么失败数据最值钱}
失败轨迹包含边界条件和恢复信息。对机器人来说，成功示范告诉模型“怎么做”，失败后修正告诉模型“哪里会错”和“如何恢复”。这类数据比重复成功演示更有边际价值。
\end{knowledgebox}

\subsection{本章小结}

机器人数据采集不是单一路径，而是多来源组合。真实部署、遥操作、仿真、合成和互联网视频各自解决不同问题。

